\begin{problem}
    Para una tarea de Estadística, una estudiante tiene que representar un conjunto de datos a través de un diagrama de cajón. Los datos que se presentan a continuación corresponden a la cantidad de anotaciones positivas que tuvieron los amigos de esta estudiante el semestre anterior.
    \[
    4 - 0 - 0 - 2 - 3 - 3 - 1 - 0 - 3
    \]

    Para construir el diagrama de cajón, ella realiza los siguientes pasos, cometiendo un error.

    \begin{itemize}
        \item Paso 1: ordena los datos de menor a mayor y determina que el dato central es \num{2}.
        \item Paso 2: identifica los valores mínimo y máximo, obteniendo \num{0} y \num{4} respectivamente.
        \item Paso 3: determina el primer cuartil como \num{1}, ya que se encuentra entre \num{0} y \num{2}, y el tercer cuartil como \num{3}, ya que se encuentra entre \num{2} y \num{4}.
        \item Paso 4: construye el diagrama de cajón, obteniendo:
    \end{itemize}

    \begin{center}
    (imagen)
    \end{center}

    ¿En cuál de los pasos la estudiante cometió el error?

    \begin{enumerate}[label=\Alph*.]
        \item En el Paso 1
        \item En el Paso 2
        \item En el Paso 3
        \item En el Paso 4
    \end{enumerate}
\end{problem}
